\documentclass[11pt]{article}
\usepackage[T1]{fontenc}
\usepackage{textcomp}
\usepackage[margin=0.75in]{geometry}
\usepackage{amsmath,amssymb,amsthm}
\usepackage{array,booktabs}
\usepackage{hyperref}
\setlength{\parindent}{0pt}
\newtheorem{theorem}{Theorem}
\theoremstyle{definition}
\newtheorem{definition}{Definition}
\title{Flow Proof Artifact: prop-48-pythagoras-converse}
\author{Generated by \texttt{flow doc proof}}
\date{Flow Proof Book}
\begin{document}
\maketitle
If the square on one side equals the sum of squares on the other two, the angle between them is right.\\[0.5em]
\textbf{Source.} euclid — Elements, Book I, Proposition 48\\[0.5em]
\bigskip
\noindent{\large \textbf{Derived fact 1.} \textit{If the square on one side equals the sum of squares on the other two, the angle between them is right} \hfill the Euclidean plane \cdot Euclid Book I \cdot Proposition 48: square sum implies a right angle}\par
\smallskip\noindent\textit{Coordinate: the Euclidean plane \cdot Euclid Book I \cdot Proposition 48: square sum implies a right angle}\par
\smallskip\noindent\textit{Source: euclid — Elements, Book I, Proposition 48}\par
\smallskip\noindent\textit{Built on: proposition 47: square on hypotenuse equals sum of squares on legs, for Book I of the Elements in on the Euclidean plane, proposition 8: side-side-side angle equality, for Book I of the Elements in on the Euclidean plane}\par
\medskip
\noindent\textbf{Goal.} If the square on one side equals the sum of squares on the other two, the angle between them is right\par
\noindent$\displaystyle \angle BAC = 90^\circ$\par
\medskip
\noindent
\begin{tabular}{@{} >{\raggedright\arraybackslash}p{0.06\textwidth} >{\raggedright\arraybackslash}p{0.47\textwidth} >{\raggedleft\arraybackslash}p{0.41\textwidth} @{}}
\textbf{\#} & \textbf{Proof} & \textbf{Mathematics} \\
\midrule
\textbf{1.} & We prove that proposition 48: square sum implies a right angle for Book I of the Elements in the Euclidean plane. & \\[0.45em]
\textbf{2.} & Let D = point\_with\_AD\_perpendicular\_to\_AB\_and\_AD\_equal\_to\_AC. & \\[0.45em]
\textbf{3.} & Let square\_on\_BD = square\_on\_side\_BD. & \\[0.45em]
\textbf{4.} & We invoke the derived fact: triangle\_ABC\_with\_BC\_squared\_equals\_AB\_squared\_plus\_AC\_squared. & \\[0.45em]
\textbf{5.} & We invoke the derived fact governing Book I of the Elements in the Euclidean plane: proposition 47: square on hypotenuse equals sum of squares on legs, for Book I of the Elements in on the Euclidean plane. & \\[0.45em]
\textbf{6.} & We invoke the derived fact governing Book I of the Elements in the Euclidean plane: proposition 8: side-side-side angle equality, for Book I of the Elements in on the Euclidean plane. & \\[0.45em]
\textbf{7.} & From step 2, step 3, step 4, step 5, and step 6, we can deduce that angle DAB equals one right angle. & $\displaystyle \angle DAB = 90^\circ$ \\[0.45em]
\textbf{8.} & We can deduce that angle BAC equals one right angle. Hence proven. & $\displaystyle \angle BAC = 90^\circ$ \\[0.45em]
\bottomrule
\end{tabular}
\label{thm:Geometry-Euclid Book I-proposition_48_square_sum_implies_a_right_angle}
\par\medskip\hrule
\end{document}
